\documentclass[11pt]{article}
\usepackage[utf8]{inputenc}
\usepackage[T1]{fontenc}
\usepackage{amsmath,amssymb}
\usepackage{booktabs}
\usepackage{hyperref}
\usepackage[margin=1in]{geometry}
\usepackage{graphicx}
\usepackage{xcolor}
\usepackage{enumitem}
\usepackage{natbib}
\bibliographystyle{plainnat}

\title{Agent Trajectory Failure Detection:\\A Benchmark for Evaluating Agent Monitoring Tools}

\author{
  Sohan Shingade \\
  Galea \\
  \texttt{sohan@galea.foo}
}

\date{May 2026}

\begin{document}
\maketitle

\begin{abstract}
As AI agent workflows become production business processes, tools that monitor these workflows must detect failures automatically---without requiring users to write evaluation rules. We introduce \textbf{Agent Trajectory Failure Detection (ATFD)}, a benchmark task that measures how well a monitoring tool can (1) detect failed agent trajectories, (2) avoid false positives on successful trajectories, and (3) correctly categorize the type of failure. We evaluate on 150 trajectories from $\tau$-bench \citep{tau2bench}, spanning retail and airline customer-service domains. We report results for Galea, a heuristic-based agent investigation system, achieving 94.4\% detection rate with 0\% false positive rate and 79\% failure-type alignment---using zero user-configured evaluation rules. We release the benchmark methodology, converter, and evaluation harness to enable comparison across agent monitoring tools.
\end{abstract}

\section{Introduction}

Agent observability tools---LangSmith, Langfuse, Arize Phoenix, Braintrust---provide trace collection and visualization for LLM-based agent workflows. However, none of these tools automatically detect whether a given agent trajectory represents a failure. Users must write custom evaluation rules, set alert thresholds, or manually review traces.

This gap matters as agent workflows increasingly handle consequential business tasks: processing refunds, modifying reservations, exchanging products, and handling customer complaints. A customer-service agent that calls the wrong tool, passes incorrect arguments, or enters a tool-call loop represents a real business failure---not a model quality issue.

Existing agent benchmarks---$\tau$-bench \citep{tau2bench}, AgentBench \citep{agentbench}, AgentBoard \citep{agentboard}, SWE-bench \citep{swebench}---evaluate \emph{agents} on task completion. No benchmark evaluates the \emph{tools that monitor agents}. We address this gap.

We propose Agent Trajectory Failure Detection (ATFD) as a benchmark task. The task is simple: given a complete agent trajectory (the full sequence of messages, tool calls, and tool results from a multi-turn interaction), determine whether the trajectory represents a failure, and if so, what kind.

\section{Task Definition}

\subsection{Input}

A complete agent trajectory consisting of:
\begin{itemize}[nosep]
  \item System prompt and user messages (customer requests)
  \item Assistant messages (agent responses)
  \item Tool calls (name, arguments, requestor)
  \item Tool results (output, error status)
  \item Termination reason
\end{itemize}

\subsection{Output}

A set of \textbf{findings}, each with:
\begin{itemize}[nosep]
  \item \textbf{Severity:} error, warning, or info
  \item \textbf{Category:} failure type (e.g., wrong action, tool loop, tool risk, policy violation)
  \item \textbf{Attribution:} which agent or tool is responsible (optional)
\end{itemize}

\subsection{Ground Truth}

Provided by $\tau$-bench's evaluation framework:
\begin{itemize}[nosep]
  \item \texttt{reward} (0.0 = complete failure, 1.0 = success)
  \item \texttt{reward\_breakdown} (per-component: DB state match, action correctness, communication completeness)
  \item \texttt{termination\_reason} (normal stop, error, timeout, max steps)
\end{itemize}

\section{Metrics}

\subsection{Detection Rate}

Percentage of failed trajectories ($\text{reward} < 1.0$) where the tool produced at least one substantive finding:
\[
\text{Detection Rate} = \frac{|\{t \in T_{\text{fail}} : \text{findings}(t) \neq \emptyset\}|}{|T_{\text{fail}}|}
\]

\subsection{False Positive Rate}

Percentage of successful trajectories ($\text{reward} = 1.0$) where the tool produced error-severity findings:
\[
\text{FP Rate} = \frac{|\{t \in T_{\text{pass}} : \text{error\_findings}(t) \neq \emptyset\}|}{|T_{\text{pass}}|}
\]

Baseline anomaly findings (cost/latency/volume deviations from project median) are excluded from the FP calculation, as these are informational and expected when the project baseline has few samples.

\subsection{Category Alignment}

Percentage of $\tau$-bench failure types correctly covered by the tool's finding categories. The mapping between $\tau$-bench failure types and tool finding categories is defined in Table~\ref{tab:category-map}.

\begin{table}[h]
\centering
\caption{Category mapping between $\tau$-bench failure types and tool finding categories.}
\label{tab:category-map}
\begin{tabular}{ll}
\toprule
$\tau$-bench failure type & Matching tool categories \\
\midrule
\texttt{wrong\_action} & tool\_error, tool\_risk\_halt, tool\_risk\_warn, loop\_suspect \\
\texttt{wrong\_state\_change} & tool\_risk\_halt, tool\_risk\_warn, failure \\
\texttt{wrong\_communication} & correctness\_no\_citation, hallucination \\
\texttt{missing\_information} & correctness\_no\_citation \\
\texttt{env\_violation} & policy\_deny, tool\_risk\_halt \\
\bottomrule
\end{tabular}
\end{table}

\subsection{Configuration Effort}

Number of evaluation rules, scorers, or alert thresholds the user must configure before the tool produces findings. Lower is better. Zero means fully automatic detection.

\section{Dataset}

We use pre-recorded simulation results from $\tau$-bench \citep{tau2bench}, a benchmark for evaluating LLM agents on customer-service tasks.

\paragraph{Domains.} Retail (100 trajectories): product exchanges, order cancellations, address modifications. Airline (50 trajectories): reservation changes, cancellations, certificate issuance.

\paragraph{Agent configuration.} GPT-4.1 (2025-04-14) with default $\tau$-bench scaffold. User simulator: GPT-4.1 (2025-04-14). Four trials per task.

\paragraph{Trajectory size.} Each trajectory contains 10--53 trace events after conversion, including user messages, model calls, tool invocations, and tool results.

\begin{table}[h]
\centering
\caption{Failure distribution across domains.}
\begin{tabular}{lccc}
\toprule
Domain & Total & Failed ($r < 1.0$) & Pass rate \\
\midrule
Retail & 100 & 22 & 78\% \\
Airline & 50 & 14 & 72\% \\
\midrule
Combined & 150 & 36 & 76\% \\
\bottomrule
\end{tabular}
\end{table}

\section{Evaluated System: Galea}

Galea is an investigation layer for agent workflows. Its heuristic investigator operates without LLM calls and detects failures through five mechanisms:

\begin{enumerate}[nosep]
  \item \textbf{Event pattern matching}---run failures, policy denials, approval rejections, tool errors
  \item \textbf{Counting}---tool-call repetition ($\geq$4 calls to the same tool flags a loop suspect)
  \item \textbf{Statistical comparison}---current trace vs.\ project baseline median ($2\times$ = warning, $5\times$ = error)
  \item \textbf{Cross-referencing}---cited documents vs.\ indexed documents in the data room
  \item \textbf{Risk scoring}---tool permission level $\times$ repetition $\times$ resource sensitivity
\end{enumerate}

No LLM calls, no embeddings, no user-configured rules. The system ingests trace events and produces findings automatically.

\paragraph{Conversion.} $\tau$-bench \texttt{SimulationRun} objects are converted to Galea's normalized \texttt{TraceEvent} format. Each message becomes the appropriate event type (\texttt{signal\_received}, \texttt{model\_called}, \texttt{tool\_called}, \texttt{tool\_completed}, \texttt{tool\_failed}). The converter is included in the open-source release.

\section{Results}

\begin{table}[h]
\centering
\caption{ATFD benchmark results for Galea (heuristic investigator, zero configuration).}
\label{tab:results}
\begin{tabular}{lccc}
\toprule
Metric & Retail ($n$=100) & Airline ($n$=50) & Combined \\
\midrule
Detection Rate & 90.9\% & 100.0\% & 94.4\% \\
False Positive Rate & 0.0\% & 0.0\% & 0.0\% \\
Category Alignment & 76.2\% & 85.0\% & 79.0\% \\
Avg.\ findings (failed) & 3.4 & 4.3 & 3.8 \\
Avg.\ findings (passed) & 2.7 & 1.7 & 2.3 \\
Config effort & 0 & 0 & 0 \\
\bottomrule
\end{tabular}
\end{table}

\subsection{Failure Analysis}

The two undetected failures in the retail domain were trajectories where the agent called the correct tools in the correct order but passed \textbf{wrong arguments} (e.g., exchanging item A for item B instead of item A for item C). The trace shape was identical to successful trajectories---only the argument values differed. Detecting these requires comparing tool arguments against expected outcomes or using LLM-backed semantic analysis, which the heuristic layer does not perform.

Galea includes an optional LLM-backed investigation layer (not evaluated here) that compares agent outputs against a mission statement derived from the input signal and project configuration. We expect this layer to close the gap on argument-level correctness failures.

\subsection{Finding Distribution}

Most common finding categories on failed trajectories:
\begin{itemize}[nosep]
  \item \texttt{tool\_risk\_warn}---state-changing tools flagged for review
  \item \texttt{loop\_suspect}---repeated tool calls indicating planning failures
  \item \texttt{anomaly\_volume}---more events than project baseline
\end{itemize}

The airline domain achieved higher category alignment (85\% vs.\ 76.2\%) because airline tasks involve more state-changing operations (reservation modifications, cancellations), which the tool-risk scorer surfaces effectively.

\section{Comparison with Existing Tools}

No existing agent observability tool provides automatic trajectory-level failure detection without user configuration.

\begin{table}[h]
\centering
\caption{Comparison of agent monitoring tools on automatic failure detection capability.}
\begin{tabular}{lcc}
\toprule
Tool & Auto-detection & Configuration required \\
\midrule
\textbf{Galea} & \textbf{Yes (heuristic)} & \textbf{None} \\
LangSmith & No & User writes eval rules \\
Langfuse & No & User defines scores \\
Arize Phoenix & No & User configures evaluators \\
Braintrust & No & User writes scorers \\
DeepEval & No & User selects metrics \\
\bottomrule
\end{tabular}
\end{table}

This benchmark is tool-agnostic. We invite submissions from any agent monitoring tool. Tools that require configuration should report what was configured (number of rules, setup time) alongside detection metrics.

\section{Reproducing Results}

The benchmark harness, converter, and evaluation code are available at:

\begin{center}
\url{https://github.com/Galea-foo/Galea/tree/main/benchmarks/tau-bench}
\end{center}

\section{Conclusion}

We introduce Agent Trajectory Failure Detection (ATFD), the first benchmark for evaluating agent monitoring tools---as opposed to agents themselves. On 150 $\tau$-bench trajectories, Galea's heuristic investigator achieves 94.4\% detection rate with 0\% false positive rate, using zero user-configured rules. We release the benchmark to enable systematic comparison across the growing ecosystem of agent observability and investigation tools.

\begin{thebibliography}{9}

\bibitem[Sierra Research(2025)]{tau2bench}
Sierra Research.
\newblock $\tau$-bench: A Benchmark for Tool-Agent-User Interaction in Real-World Domains.
\newblock \url{https://github.com/sierra-research/tau2-bench}, 2025.

\bibitem[Liu et~al.(2024)]{agentbench}
Xiao Liu, Hao Yu, Hanchen Zhang, et~al.
\newblock AgentBench: Evaluating LLMs as Agents.
\newblock In \emph{ICLR}, 2024.

\bibitem[Ma et~al.(2024)]{agentboard}
Chang Ma, Junlei Zhang, Zhihao Zhu, et~al.
\newblock AgentBoard: An Analytical Evaluation Board of Multi-turn LLM Agents.
\newblock In \emph{NeurIPS}, 2024.

\bibitem[Jimenez et~al.(2024)]{swebench}
Carlos~E. Jimenez, John Yang, Alexander Wettig, et~al.
\newblock SWE-bench: Can Language Models Resolve Real-World GitHub Issues?
\newblock In \emph{ICLR}, 2024.

\bibitem[Yao et~al.(2025)]{beyondblackbox}
Yifan Yao, Yiran Wu, et~al.
\newblock Beyond Black-Box Benchmarking: Observability, Analytics, and Optimization of Agentic Systems.
\newblock In \emph{KDD}, 2025.

\end{thebibliography}

\end{document}
